\documentclass[../DoAn.tex]{subfiles}
\begin{document}

\begin{center}
    \Large{\textbf{TÓM TẮT NỘI DUNG ĐỒ ÁN}}\\
\end{center}
\vspace{1cm}
Trong bối cảnh chuyển đổi số ngành giáo dục ngày càng được đẩy mạnh, nhu cầu về phát triển học tập với các thiết bị, hệ thống hỗ trợ học tập thông minh ngày càng trở nên cần thiết. Đối với các học phần đại học, lượng học phần, quy chế về các học phần, các câu hỏi thường gặp thường được lưu trữ ở nhiều nơi khiến cho việc tìm kiếm mất rất nhiều thời gian. Nắm bắt được điều này, em đã quyết đình làm đồ án về đề tài "Trợ lý ảo hỏi đáp về học phần của trường Đại Học Bách Khoa Hà Nội" sẽ giúp sinh viên phần nào việc tra cứu trở nên dễ dàng hơn.\\
Trong đồ án lần này, với sự phát triển của các mô hình ngôn ngữ lớn LLM. Việc xây dựng các hệ thống hỏi đáp đã đạt được nhiều bước tiến vượt bậc. Vì vậy, em đã lựa chọn hướng tiếp cần RAG để giúp mô hình khai tác được nguồn dữ liệu của trường, giảm hiện tượng tạo ra các thông tin không chính xác và cho phép cập nhật dữ liệu theo giời gian mà không cần huấn luyện lại mô hình. Hệ thống được xây dựng với quy trình thu thập dữ liệu về thông tin các học phầ, tạo cơ sở dữ liệu vector, truy xuất các văn bản liên quan và sử dụng LLM để trả lời câu hỏi dựa trên nội dung tương ứng.\\
Đồ án tập trung vào vào việc nghiên cứu, xây dựng và triển khai thử nghiệm một trợ lý ảo hỏi đáp về thông tin học phần giúp hỗ trợ sinh viên tra cứu các thông tin liên quan một cách dễ dàng và hiệu quả hơn. Nâng cao trải nghiệm của sinh viên trong môi trường đại học


\begin{flushright}
Sinh viên thực hiện\\
\begin{tabular}{@{}c@{}}
\textit{(Ký và ghi rõ họ tên)} \\ \\ \\
Phạm Thành Huy
\end{tabular}
\end{flushright}

\end{document}